
\begin{dang}{CON LẮC ĐƠN}
%$\bullet$ \textbf{Các công thức cơ bản}
%\begin{bang}[tabularx={X|X|X|X}]{}
%Chu kì&Tần số&Tần số góc&Chuyển đổi\\\hline
%\[T=2\pi\sqrt{\dfrac{\ell}{g}}\]&\[f=\dfrac{1}{2\pi}\sqrt{\dfrac{g}{\ell}}\]&\[\omega=\sqrt{\dfrac{g}{\ell}}\]&\[s=\ell.\alpha;\  S_0=\ell.\alpha_0\]\\\hline
%\end{bang}
%\begin{note}
%Ta dùng mối quan hệ tỉ lệ để rút ra mối liên hệ chu kì hoặc tần số trong các trường hợp chiều dài dây treo thay đổi.
%\end{note}
%$\bullet$ \textbf{Thay đổi chu kì do tác động bên ngoài}
%\begin{bang}[tabularx={X|X|X}]{}
%Do nhiệt độ&Đưa lên cao&Đưa lên thiên thể\\\hline
%\[\ell'=\ell_0\left(1+\alpha\Delta t\right)\]&\[g'=\dfrac{GM}{\left(R+h\right)^2}\]&\[g'=\dfrac{GM'}{R'^2}\]\\\hline
%\end{bang}}
%\begin{itemize}
%\item Gọi T,  T' lần lượt là chu kì của đồng đồ đúng và chu kì của đồng hồ sai.
%\item Giả sử hai đồng hồ bắt đầu hoạt động cùng một lúc và đến một thời điểm số chỉ của chúng lần lượt là t và t'.
%\item Theo nguyên  tắc cấu tạo đồng hồ quả lắc thì: $t_{\text{đúng}}T_{\text{đúng}}=t_{\text{sai}}T_{\text{sai}}$.
%\end{itemize}
\end{dang}
